\section*{Problemas}

\subsection*{Enunciados de los problemas}

% Recordar
\begin{problema}
	\label{prob:06f83cc}
	Un centro de atención telefónica recibe llamadas de manera independiente a una tasa promedio constante $\lambda$ por minuto. Enuncie, sin realizar ningún cálculo, la función de masa de probabilidad $P(X=k)$ de $X\sim\text{Pois}(\lambda)$ y su soporte $\mathcal{S}_X$.
\end{problema}

% Comprender
\begin{problema}
	\label{prob:c35cfa7}
	Un sistema industrial registra fallas de una máquina con tasa $0.5$ fallas/hora, bajo un proceso de Poisson. Explique por qué el número de fallas en un turno de $4$ horas sigue $\text{Pois}(4\times0.5)=\text{Pois}(2)$, y por qué modelar el mismo turno como $n=4$ ensayos Bernoulli horarios independientes con $p=1-e^{-0.5}$ produce exactamente la misma probabilidad de al menos una falla, $P(X\ge1)=1-e^{-2}$.
\end{problema}

% Aplicar
\begin{problema}
	\label{prob:8fd3390}
	El número de personas que llegan por día a una sala de urgencias sigue $\text{Pois}(\lambda=5)$.
	\begin{enumerate}
		\item Calcule $P(X\le3)$, la probabilidad de que lleguen a lo sumo $3$ personas en un día.
		\item Calcule $P(X\ge8)$, la probabilidad de que lleguen al menos $8$ personas.
	\end{enumerate}
\end{problema}

% Analizar
\begin{problema}
	\label{prob:8dad711}
	Sean $X\sim\text{Pois}(\lambda_1)$ y $Y\sim\text{Pois}(\lambda_2)$ independientes (puede usar sin demostrar que $X+Y\sim\text{Pois}(\lambda_1+\lambda_2)$). Demuestre que la distribución condicional de $X$ dado $X+Y=n$ es Binomial con parámetros $n$ y $p=\lambda_1/(\lambda_1+\lambda_2)$, y aplique el resultado con $\lambda_1=3$, $\lambda_2=5$, $n=8$ para calcular $P(X=3\mid X+Y=8)$.
\end{problema}

% Evaluar
\begin{problema}
	\label{prob:5e9408a}
	Sean $X_1,\dots,X_{100}$ i.i.d. $\text{Pois}(\lambda=2)$, con $S_{100}=\sum X_i \sim \text{Pois}(200)$. Por el Teorema Central del Límite, $S_{100}$ se aproxima por $N(200,200)$. Evalúe la calidad de esta aproximación calculando $P(S_{100}\le220)$ de forma exacta (vía la Poisson de suma) y mediante la aproximación Normal con corrección de continuidad, y determine si el error absoluto es lo suficientemente pequeño para justificar el uso práctico de la aproximación.
\end{problema}

% Crear
\begin{problema}
	\label{prob:9dc367e}
	Diseñe un ejemplo original (contexto y tasa $\lambda$) que se modele naturalmente con una distribución de Poisson, distinto de los ejemplos de centros de llamadas, manufactura o servidores web ya vistos en esta sección. Plantee la pregunta de probabilidad de interés y calcule $\mathbb{E}[X]$ y $\text{Var}(X)$ para el parámetro que eligió.
\end{problema}

\subsection*{Soluciones detalladas}

% Recordar
\begin{solproblema}
	$P(X=k) = \dfrac{e^{-\lambda}\lambda^k}{k!}$, para $k \in \mathcal{S}_X = \set{0,1,2,\dots}$.
\end{solproblema}

% Comprender
\begin{solproblema}
	En un proceso de Poisson, el número de eventos en un intervalo de tiempo escala linealmente con su duración: si la tasa es $0.5$ fallas/hora, en $4$ horas la tasa esperada es $4\times0.5=2$, de modo que el conteo sigue $\text{Pois}(2)$, y $P(X\ge1)=1-e^{-2}\approx0.8647$. Al discretizar el mismo turno en $4$ ensayos Bernoulli horarios independientes, la probabilidad de al menos una falla en una hora aislada es $p=1-e^{-0.5}\approx0.3935$ (el complemento de $P(\text{Pois}(0.5)=0)$), y la probabilidad de al menos una falla en las $4$ horas es $1-(1-p)^4 = 1-(e^{-0.5})^4 = 1-e^{-2}$, idéntica a la anterior. Ambos modelos coinciden porque son dos formas equivalentes de contar el mismo proceso de llegadas raras: uno en tiempo continuo (Poisson) y otro discretizado en subintervalos (Bernoulli/Binomial), y ambos comparten la misma probabilidad exacta de "cero eventos" en cada subintervalo.
\end{solproblema}

% Aplicar
\begin{solproblema}
	Con $\lambda=5$:
	\begin{enumerate}
		\item $P(X\le3) = e^{-5}\left(1+5+\dfrac{25}{2}+\dfrac{125}{6}\right) \approx 0.2650$.
		\item $P(X\ge8) = 1-P(X\le7) \approx 1-0.8670 = 0.1330$.
	\end{enumerate}
\end{solproblema}

% Analizar
\begin{solproblema}
	Por definición de probabilidad condicional e independencia:
	\begin{align*}
		P(X=k\mid X+Y=n) = \dfrac{P(X=k)P(Y=n-k)}{P(X+Y=n)} = \dfrac{\dfrac{e^{-\lambda_1}\lambda_1^k}{k!}\cdot\dfrac{e^{-\lambda_2}\lambda_2^{n-k}}{(n-k)!}}{\dfrac{e^{-(\lambda_1+\lambda_2)}(\lambda_1+\lambda_2)^n}{n!}} = \comb{n}{k}\left(\dfrac{\lambda_1}{\lambda_1+\lambda_2}\right)^k\left(\dfrac{\lambda_2}{\lambda_1+\lambda_2}\right)^{n-k},
	\end{align*}
	que es la PMF de $\text{Bin}(n,p)$ con $p=\lambda_1/(\lambda_1+\lambda_2)$. Con $\lambda_1=3,\lambda_2=5,p=3/8$: $P(X=3\mid X+Y=8) = \comb{8}{3}\left(\dfrac{3}{8}\right)^3\left(\dfrac{5}{8}\right)^5 \approx 0.2815$.
\end{solproblema}

% Evaluar
\begin{solproblema}
	Exacto: $P(S_{100}\le220) = \sum_{k=0}^{220}\dfrac{e^{-200}200^k}{k!} \approx 0.9862$ (valor de referencia bajo $\text{Pois}(200)$). Aproximación Normal con corrección de continuidad: $P(S_{100}\le220) \approx \Phi\!\left(\dfrac{220.5-200}{\sqrt{200}}\right) = \Phi(1.4496) \approx 0.9265$. El error absoluto entre ambos métodos es considerable ($\approx0.06$), mayor de lo que se esperaría para $\lambda=200$; esto indica que, aunque $n=100$ es grande, la aproximación Normal en la cola (donde la probabilidad exacta ya es alta) puede ser menos precisa que en el centro de la distribución, y conviene usar la Poisson exacta (o software estadístico) en vez de la aproximación Normal cuando se requiere precisión en las colas.
\end{solproblema}

% Crear
\begin{solproblema}
	\textbf{Ejemplo:} una biblioteca registra en promedio $\lambda=4$ solicitudes de préstamo interbibliotecario por día, modeladas como un proceso de Poisson. Pregunta de interés: ¿cuál es la probabilidad de que en un día dado se reciban más de $6$ solicitudes, $P(X>6)$? Los momentos son $\mathbb{E}[X]=\text{Var}(X)=\lambda=4$. La probabilidad se calcularía como $P(X>6) = 1-\sum_{k=0}^{6}\dfrac{e^{-4}4^k}{k!}$.
\end{solproblema}
